\chapter{Limitations and Future Work}
\label{chap:limitations-future-work}

\section{Current prototype limitations}

The final implementation is a research prototype. It is substantially more complete and
more measurable than the initial implementation, but it should not be presented as a
production-ready fair-exchange system. The prototype runs on a local development stack,
uses local accounts, stores artifacts through a local backend, and measures gas on a local
Hardhat network. These choices are appropriate for a semester project focused on protocol
variants and gas measurements, but they leave several engineering and security tasks open.

The first limitation is the user model. The application uses predefined development
addresses and role selection through the local interface. It does not implement a
production authentication flow where the vendor, buyer, sponsor, and dispute sponsors
connect through real wallets and explicitly sign every role-dependent action. Some
authorization checks are implemented on-chain, and the buyer-side dispute sponsor
authorization is bound to the chain id and contract address. However, the surrounding user
management remains a proof-of-concept environment.

The second limitation is the backend model. The implementation stores ciphertext artifacts,
precontract data, and local metadata through the application backend and SQLite database.
The on-chain commitments make tampering detectable, but the backend is still an
availability dependency. If a production deployment wants to avoid this dependency, it
would need a more robust artifact distribution layer, for example content-addressed
storage, replication, or a clearly specified off-chain availability protocol.

The third limitation is the separation between the optimized paths and the inherited
account-abstraction path. The normal direct contract keeps the ERC-4337-compatible behavior
of the initial implementation. The clone-based optimized variants are measured through
direct calls and are intentionally thinner. This is the right choice for gas optimization,
but it means that the project does not yet provide one unified production interface where
all optimized variants are available through the same wallet and bundler flow.

The fourth limitation is proof coverage. The test suite covers the main targeted scenarios:
normal optimistic execution, the three main optimistic variants, self-sponsored disputes,
hardcoded \(\mathsf{SHA256}\), same-scope gas comparisons, and the 1 GiB hardcoded
benchmark. This is strong evidence for the implemented paths, but it is not a formal proof
of the Solidity implementation. A production version would require fuzzing, invariant
tests, adversarial tests around malformed proofs, and preferably a formal specification of
the state machines and commitment formats.

\section{Verification workflow limitations}

The implementation contains verification mechanisms, but they are not yet a complete
production verification workflow. This distinction is important. The buyer-facing
precontract modal exposes a ``Verify commitment'' action that calls
\path{/api/precontracts/verify}. That route invokes the native
\texttt{check\_precontract\_cli} binary and recomputes the public commitment data from the
stored ciphertext, description, commitment, and opening value. Thus, the codebase does
include a concrete verification path for the received precontract data.

However, this verification is not enforced as a backend precondition for acceptance. The
acceptance route \path{/api/precontracts/accept} currently sets the contract as accepted
from the submitted identifier. It does not require a previously stored successful
verification result, nor does it bind the acceptance transaction to a verified
\(h_{\mathit{ct}}\) and \(h_{\mathit{circuit}}\) record. In the proof-of-concept UI, the
buyer is expected to click the verification button before accepting. In a production
version, acceptance should be blocked unless verification succeeds and the verified values
are persisted and bound to the exchange.

The same limitation exists in the dispute workflow. The contracts verify the submitted
commitments and witnesses on-chain during Step 8, and the native benchmarks generate the
data needed for complete measured disputes. Nevertheless, the web API for proof generation
is still partial: \path{/api/proofs/compute} currently supports only state 4
(\texttt{WaitVendorDataRight}). It is therefore not a complete production API for every
possible Step 8 state. A full implementation should support all Step 8 cases used by the
protocol and should make the actor-side local verification explicit before the actor signs
the next dispute action.

This also affects runtime interpretation. The 1 GiB measurements include native generation
of buyer-side \(hpre_i\) values and the on-chain \texttt{giveOpinion} calls, but they do
not isolate a complete production UI workflow in which \(V\) locally recomputes every
buyer response before deciding what to sign. For conservative protocol-level performance
tables, such verification time should be counted for the verifying actor, even if the
current benchmark does not expose it as a separate UI measurement.

\section[Browser and WASM limits]{Browser and WebAssembly limitations for large files}

One concrete outcome of the project is that the browser/WASM path is not the right path for
very large files. The WebAssembly integration remains useful for the interface and for
small or medium examples, but the large-file benchmarks rely on the native Rust CLI. This
is not accidental. A 1 GiB exchange requires reading large files, encrypting them,
constructing commitments, and generating dispute witnesses. Doing this inside a browser or
browser-adjacent JavaScript runtime creates avoidable memory pressure and makes results
less reliable.

The measured 1 GiB path should therefore be described precisely: it is a native CLI
benchmark integrated with the Solidity dispute contracts, not a complete browser-based
1 GiB user experience. This is still an important result, because it demonstrates that the
protocol artifacts and the on-chain dispute can be produced and checked at this scale. It
also clarifies the architecture that a future large-file application should use.

A natural future direction is to make the native path a first-class backend worker. Instead
of asking the frontend to handle large cryptographic computations directly, the interface
could delegate precontract generation, buyer-side \(hpre_i\) computation, and dispute
witness generation to a native service. The frontend would then orchestrate the exchange
and display progress, while the native worker would handle file streaming and cryptographic
artifact generation.

This future work should also include streaming and incremental verification. The current
large-file benchmark proves feasibility, but it is still benchmark-oriented. A production
implementation would need careful progress reporting, resumability, cancellation, disk
space management, and checks that intermediate artifacts are not corrupted or mismatched
with the on-chain commitments.

\section[Registry-first architecture]{Toward a registry or factory-first architecture}

The final implementation already introduces a factory for clone-based optimistic variants,
but it does not fully redesign the protocol around a registry. Each exchange still gets a
dedicated optimistic account or clone, and disputes still deploy dispute-specific accounts.
This is a strong improvement over redeploying full optimistic bytecode every time, but it
is not necessarily the end point of the architecture.

A more radical future architecture would use a factory-first or registry-first model. In
such a design, reusable contracts would be deployed once, and each exchange would be
represented by an \texttt{exchangeId} in a shared registry or by a minimal deterministic
proxy. The objective would be to reduce repeated deployment costs even further and make
Step 1+2 fusion easier to express. This would be a larger redesign than the one carried out
in this project, because state layout, role authorization, dispute deployment, and artifact
references would all need to be specified around the registry.

The main advantage would be gas efficiency and clearer lifecycle management. A registry can
avoid deploying large bytecode for every exchange and can make it easier to reuse a sponsor
approval if a buyer abandons the flow before paying. It could also centralize event
emission, indexing, and off-chain discovery of exchanges.

The main risk is complexity. A shared registry must avoid cross-exchange state confusion,
must isolate deposits correctly, and must keep the dispute logic auditable. The current
clone architecture is a conservative intermediate step: it captures most of the deployment
gain while keeping each exchange isolated. Moving to a registry should therefore be treated
as a separate design task, not as a small patch.

\section[Step 1+2 fusion for S not equal to B]{Step 1+2 fusion for \(S \neq B\)}

The project implemented the trivial Step 1+2 fusion for \(S=B\): the buyer is also the
initial sponsor, so deployment and buyer payment can be combined. A more difficult and
interesting future direction is fusing Step 1+2 when \(S \neq B\). The goal is to avoid a
situation where \(S\) deploys or initializes a contract and then loses that deployment if
\(B\) never pays.

One possible design is to move the loading of the precontract parameters from Step 1 to
Step 2. The sponsor \(S\) would sign an approval for a specific set of precontract
parameters. Then \(B\) would submit, in one transaction, the buyer payment, the precontract
parameters, and the signature from \(S\). If \(B\) does not submit this transaction, the
sponsor has not locked a dedicated contract instance for that buyer. The same idea could be
combined with deterministic factory deployment or with a registry entry created only when
the buyer arrives with a valid approval.

This direction has to be studied carefully with ERC-4337. If a bundler or paymaster is
expected to sponsor the transaction, it must be able to validate the UserOperation before
execution. That means the signature from \(S\), the expected buyer payment, the selected
precontract parameters, and the target deterministic address or registry entry must all be
checkable from the UserOperation data. If this validation requires executing state changes
that are only available after deployment, the design becomes incompatible with a clean
pre-sponsorship check.

For this reason, Step 1+2 fusion for \(S \neq B\) is best treated as future work on the
deployment architecture. It is closely connected to the registry/factory-first direction.
The current implementation provides the measurements and specialized contracts needed to
motivate this work, but it does not implement the full signed-precontract deployment
protocol.

\section{Cleaner circuit encoding}

The inherited V2 circuit encoding is functional for the generated circuits used in the
benchmarks, but it is not the cleanest format for a final protocol specification. Each gate
is encoded on 64 bytes with an opcode, son indices, parameters, and zero padding. Some gate
families have variable arity. In particular, SHA2 can use one or two inputs, and CONST can
use zero or one input. In the current format, arity is not encoded as an explicit field or
as a distinct opcode family.

This creates avoidable ambiguity at the specification level. Even if the current generated
circuits and parser checks behave consistently for the tested paths, a clean format should
not rely on inferring arity from padding conventions or from the number of witness values
provided in Step 8. A paper-ready and audit-ready version should make the interpretation of
every 64-byte gate canonical.

A future V3 circuit format should therefore introduce either explicit arity or distinct
opcodes. For example, it could separate:

\begin{itemize}
  \item the first SHA-256 compression gate from subsequent SHA-256 compression gates;
  \item CONST-without-input from CONST-with-input;
  \item any future gate family whose interpretation depends on optional sons.
\end{itemize}

This change would affect the Rust circuit generator, the Solidity evaluator, the
hardcoded \(\mathsf{SHA256}\) reconstruction library, the test vectors, and the computation
of \(h_{\mathit{circuit}}\). It was intentionally not patched during the gas-optimization
phase because it would have changed the commitment format and made comparison with the
initial implementation less direct. It is nevertheless one of the most important future
security cleanups.

\section{Merkle domain separation}

The accumulator construction used for \(h_{\mathit{circuit}}\), \(h_{\mathit{ct}}\), and
\(hpre\) should also be revised in a future version. The current implementation uses
\(\mathrm{keccak256}\) for leaves and internal nodes without an explicit domain separator,
and lonely nodes at an odd level can be copied upward. This convention is inherited from
the existing implementation and is mirrored by the Solidity verifier for compatibility with
the generated proofs.

The limitation is structural. If leaves and internal nodes use the same hash domain, a
64-byte encoded leaf has the same input length as the concatenation of two 32-byte child
hashes. This creates a classical ambiguity risk in Merkle-style constructions. Similarly,
copying a lonely node upward can make different tree shapes share the same root unless the
proof format completely rules out the ambiguity.

A cleaner accumulator version should domain-separate leaves and internal nodes. A simple
candidate is:

\[
\begin{aligned}
  \mathsf{leaf}_i =
  \mathrm{keccak256}(0x00 \parallel \mathsf{encodedLeaf}_i), \\
  \mathsf{node} =
  \mathrm{keccak256}(0x01 \parallel \mathsf{left} \parallel \mathsf{right}).
\end{aligned}
\]

It should also replace copy-up behavior with an explicit single-child rule, for example a
separate prefix for promoted nodes. This change must be applied consistently to all
accumulators and proofs. Like the circuit encoding change, it would create a new commitment
version. It is therefore future work rather than a safe last-minute patch to the measured
implementation.

\section{Wallet integration}

The current interface is a development interface. It allows the project to demonstrate the
protocol flow, select roles, create precontracts, accept exchanges, and trigger contract
actions. It does not yet provide a production wallet experience.

A production wallet integration should replace local role selection by explicit wallet
authentication. Each actor should connect through a wallet, see the exact action being
signed, and sign the relevant message or transaction. For off-chain authorizations, the
signed payload should include the chain id, contract address or deterministic deployment
parameters, exchange identifier, role, amount, deadline, and any commitment being approved.
This would reduce ambiguity and make signatures auditable by the user.

Wallet integration also matters for the dispute path. During a dispute, the buyer and
vendor must understand whether they are submitting a challenge response, giving an opinion,
or providing Step 8 witnesses. A production interface should expose the protocol state in
human-readable form and prevent users from signing stale or wrong-state actions.

This work is mostly engineering, but it is security-relevant. A fair-exchange protocol can
be correct on-chain and still be unsafe in practice if users cannot reliably understand
what they are authorizing. The current implementation is sufficient for experiments and
demonstrations; a deployed system would require this layer.

\section{Production deployment considerations}

A production deployment would require several additional checks beyond the scope of the
semester project. First, the contracts should be audited with the final commitment format
fixed. The current code was tested extensively for the targeted scenarios, but an audit
would need to review all state transitions, payout paths, timeout behavior, malformed proof
inputs, and edge cases around deposits and sponsor roles.

Second, the deployment model should be finalized before optimization continues. The final
project implementation already shows that deploying one large contract per exchange is not
ideal. Future work should decide whether the long-term architecture is clone-based,
registry-based, or a hybrid. That decision affects gas, upgrade strategy, indexing,
precontract signatures, and ERC-4337 compatibility.

Third, the timeout and fee parameters should be calibrated for the target chain. Local
Hardhat measurements do not capture public-network congestion, fee volatility, block time
variance, or MEV/mempool behavior. A production system would need conservative timeouts and
clear rules for who bears gas risk in the optimistic path and in dispute.

Fourth, artifact availability must be specified. The protocol can detect tampering through
commitments, but it still needs the ciphertext and witness data to be available when an
actor wants to verify or dispute an exchange. A production deployment should therefore
define where artifacts are stored, how they are addressed, how they are replicated, and how
users recover from backend unavailability.

Finally, the next gas targets are now well identified. The optimistic path has already been
reduced substantially through clones. The remaining dominant on-chain costs are dispute
deployment, \texttt{triggerDispute}, and Step 8 evaluation, especially for AES-related gates
and the associated \(h_{\mathit{ct}}\)/\(hpre\) proofs. Future optimization should focus on
these components rather than reintroducing a monolithic contract with more flags.

In summary, the project reached the intended research-prototype milestone: it implemented
and measured the main variants, corrected the architecture after the monolithic negative
result, and demonstrated a 1 GiB hardcoded dispute path. The next version should not simply
add more patches. It should clean the cryptographic formats, settle the deployment
architecture, integrate real wallets, and then optimize the remaining dispute bottlenecks
on top of that cleaner foundation.
